% Chapter 2

\chapter{文献综述}

\section{预期协同度量}
由于信息、市场认知、情绪等诸多因素的影响，
投资者对于不同证券存在预期差异\cite{barber2009retail}\cite{hong2000bad}。
预期差异使得不同投资者对同一证券形成买卖意愿，进而推动市场交易和证券定价，
是市场流动性和价格发现机制的重要来源\cite{hong2007disagreement}。

然而，在传统的金融实证研究中，预期协同度的度量较为困难。
传统的预期协同度量较为困难，主要集中于调研的宏观预期指标、代理指标、分析师预期分歧和文本分析\cite{tang2025macro}。

宏观预期指标通过调研消费者信心来反映市场预期，如密歇根大学消费者信心指数。
Lemmon等学者使用该指数作为投资者乐观程度的代理变量\cite{lemmon2004consumer}，分析情绪对于小市值股票溢价的影响。
然而，这种指标构建往往反映宏观经济预期，对于个体投资者的预期协同度量有限。

代理指标


\section{预期协同与资本市场定价}



\section{RL在投资组合管理中的应用}